\section{Medidas de Dispersão}

Além de ver onde nossos dados estão centrados, é desejável saber o quanto eles variam,
geralmente entorno do próprio centro.
Para esse fim, utiliza-se as medidas de dispersão.

Tomando a média como centro, podemos obter, como medida de dispersão, o desvio médio
absoluto, que corresponde a uma estimativa do erro médio cometido ao substituir todos
os dados pela prórpria média, indicando, também, se essa é uma boa medida de resumo dos
dados.
\[
	\dm{X} = \frac{1}{n} \sum_{i=1}^n |x_i - \bar x| \quad e \quad
\]
Outra candidata a medida de dispersão é a variância, que, como visto na seção ??,
corresponde ao segundo momento central de uma distribuição.
Numa amostra, ela pode ser interpretada como o erro quadrático médio ao substituir todos
os dados pela média, sendo obtida, geralmente para amostras grandes, pela fórmula
\[
	\var{X} = \frac{1}{n} \sum_{i=1}^n |x_i - \bar x|^2.
\]
Ao elevar as diferenças ao quadrado, desvios maiores tem mais peso do que os menores.
Como a variância possui unidade igual ao quadrado da unidade do atributo original,
então sua interpretação pode ser prejudicada, e, com isso, o desvio-padrão,
dado como
\[
	\dep{X} = \sqrt{\var{X}},
\]
pode ser usado para substituir a variância caso deseje-se que o erro esteja na unidade
dos dados.

O desvio absoluto médio e o desvio-padrão possuem interpretações e aplicações semelhantes
ao avaliar a dispersão entorno da média.
No entanto, o primeiro nos oferece uma noção mais intuitiva de dispersão, pois todos os
desvios possuem o mesmo peso independente da magnitude.
Já o desvio-padrão, uma vez que atribui pesos distintos a erros conforme a sua magnitude,
é preferível quando objetiva-se detectar rapidamente a presença de realizações muito distantes
da média, o que pode indicar valores atípicos.
